% Part: incompleteness
% Chapter: representability-in-q
% Section: undecidability

\documentclass[../../../include/open-logic-section]{subfiles}

\begin{document}

\olfileid[pt]{inc}{req}{und}
\olsection{Indecidibilidade}

Chamamos uma teoria $\Th{T}$ de \emph{indecidível} se não existe um
procedimento computacional que, depois de um número finito de passos e
sem falhar, forneça uma resposta correta à pergunta ``$\Th{T}$
prova~$!A$?'', para qualquer sentença~$!A$ na linguagem de~$\Th{T}$.
Assim, $\Th{Q}$ seria decidível sse existisse um procedimento
computacional que decidisse, dada uma sentença~$!A$ na linguagem da
aritmética, se $\Th{Q} \Proves !A$ ou não. Podemos tornar isso mais
preciso perguntando: a relação~$\Prov[\Th{Q}](y)$, que vale de~$y$ sse
$y$ é o número de G\"odel de uma sentença demonstrável em~$\Th{Q}$, é
recursiva? A resposta é: não.

\begin{thm}
$\Th{Q}$ é indecidível, O.s.\@, a relação
\[
\Prov[\Th{Q}](y) \defiff \fn{Sent}(y) \land
\lexists[x][\Prf[\Th{Q}](x, y)]
\]
não é recursiva.
\end{thm}

\begin{proof}
Suponha que fosse. Então poderíamos resolver o problema da parada da
seguinte maneira. Dados $e$ e $n$, sabemos que
$\cfind{e}(n) \fdefined$ sse existe um~$s$ tal que $T(e, n, s)$, em que
$T$ é o predicado de Kleene de
\olref[cmp][rec][nft]{thm:kleene-nf}. Como $T$ é recursivo primitivo,
ele é representável em~$\Th{Q}$ por uma fórmula $!B_T$, isto é,
$\Th{Q} \Proves !B_T(\num{e}, \num{n}, \num{s})$ sse $T(e, n, s)$.
Se $\Th{Q} \Proves !B_T(\num{e}, \num{n}, \num{s})$, então também
$ \Th{Q} \Proves \lexists[y][!B_T(\num{e}, \num{n}, y)]$. Se não
existe tal $s$, então $\Th{Q} \Proves \lnot
!B_T(\num{e}, \num{n}, \num{s})$ para todo~$s$. Mas $\Th{Q}$ é
$\omega$-consistente, O.s.\@, se
$\Th{Q} \Proves \lnot !A(\num{n})$ para todo~$n \in \Nat$, então
$\Th{Q} \Proves/ \lexists[y][!A(y)]$. Sabemos disso porque os axiomas
de $\Th{Q}$ são verdadeiros no modelo padrão~$\Struct{N}$. Logo,
$\Th{Q} \Proves/ \lexists[y][!B_T(\num{e}, \num{n}, y)]$. Em outras
palavras, $\Th{Q} \Proves \lexists[y][!B_T(\num{e}, \num{n}, y)]$ sse
existe um $s$ tal que $T(e, n, s)$, O.s.\@, sse
$\cfind{e}(n) \fdefined$. A partir de $e$ e~$n$, podemos computar
$\Gn{\lexists[y][!B_T(\num{e}, \num{n}, y)]}$; seja $g(e, n)$ a função
recursiva primitiva que faz isso. Assim,
\[
h(e, n) =
\begin{cases}
1 & \text{se $\Prov[\Th{Q}](g(e, n))$}\\
0 & \text{caso contrário}.
\end{cases}
\]
Isso mostraria que $h$ é recursiva se $\Prov[\Th{Q}]$ fosse. Mas~$h$
não é recursiva, por \olref[cmp][rec][hlt]{thm:halting-problem}; logo,
$\Prov[\Th{Q}]$ também não pode ser.
\end{proof}

\begin{cor}
A lógica de primeira ordem é indecidível.
\end{cor}

\begin{proof}
Se a lógica de primeira ordem fosse decidível, a demonstrabilidade
em~$\Th{Q}$ também seria, pois $\Th{Q} \Proves !A$ sse
$\Proves !T \lif !A$, em que $!T$ é a conjunção dos axiomas de~$\Th{Q}$.
\end{proof}

\end{document}
